% =====================================================================
%  Цель и задание лабораторной работы.
% =====================================================================

\textbf{Цель:} изучить и отработать практические навыки применения
методов автоматического распознавания языка текстовых документов.

\vspace{1em}

\textbf{Задание (вариант 25)}: разработать систему автоматического
распознавания языка текста со следующей комбинацией признаков варианта:
\begin{itemize}
  \item \textbf{Языки текста} -- французский и английский;
  \item \textbf{Формат документа} -- HTML;
  \item \textbf{Реализуемые методы} -- метод частотных слов, алфавитный
  метод и нейросетевой метод.
\end{itemize}

\vspace{1em}

Согласно методическим указаниям, для каждого метода необходимо:
\begin{enumerate}
  \item построить тренировочный набор текстов для распознаваемых языков;
  \item в соответствии с методом построить поисковый образ языка (ПОЯ)
  на основании тренировочного набора;
  \item для каждого входного документа построить поисковый образ
  документа (ПОД) по тем же правилам;
  \item рассчитать метрику расстояния между ПОД и каждым ПОЯ и выбрать
  язык с минимальным значением метрики в качестве результата;
  \item сравнить результаты распознавания разными методами по точности
  и затраченному времени.
\end{enumerate}

Кроме того, система должна:
\begin{itemize}
  \item на выходе выдавать активную ссылку на документ, результат
  идентификации языка для отдельного текста и сводную статистику по
  всей тестовой коллекции;
  \item предоставлять средства сохранения результата в файл и печати;
  \item иметь простой и доступный пользовательский интерфейс со
  справочными подсказками.
\end{itemize}
